\documentclass[10pt]{article}

\usepackage[a4paper,margin=1.7cm]{geometry}
\usepackage[T1]{fontenc}
\usepackage{mathpazo}
\usepackage{amsmath}
\usepackage{microtype}
\usepackage{booktabs}
\usepackage{tabularx}
\usepackage{array}
\usepackage{graphicx}
\usepackage{float}
\usepackage{xcolor}
\usepackage{titlesec}
\usepackage{caption}
\usepackage{enumitem}
\usepackage{fancyhdr}
\usepackage{hyperref}
\usepackage{xurl}

\definecolor{navy}{HTML}{17324D}
\definecolor{blue}{HTML}{2A6F97}
\definecolor{teal}{HTML}{2A9D8F}
\definecolor{purple}{HTML}{7B2CBF}
\definecolor{orange}{HTML}{D97706}
\definecolor{green}{HTML}{198754}
\definecolor{red}{HTML}{B42318}
\definecolor{muted}{HTML}{627D98}

\hypersetup{
  colorlinks=true,
  linkcolor=blue,
  urlcolor=purple,
  pdftitle={Higgs-small cross-model FeatureGraph transfer},
  pdfauthor={GigaEvo}
}
\titleformat{\section}{\Large\bfseries\color{navy}}{\thesection}{0.65em}{}
\titleformat{\subsection}{\large\bfseries\color{blue}}{\thesubsection}{0.55em}{}
\captionsetup{font=small,labelfont={bf,color=navy}}
\setlength{\parindent}{0pt}
\setlength{\parskip}{5pt}
\setlength{\tabcolsep}{5pt}
\renewcommand{\arraystretch}{1.14}
\emergencystretch=3em
\setlist[itemize]{leftmargin=1.4em,itemsep=2pt,topsep=3pt}

\pagestyle{fancy}
\fancyhf{}
\fancyhead[L]{\small\color{muted}Higgs-small FeatureGraph transfer}
\fancyhead[R]{\small\color{muted}23 July 2026}
\fancyfoot[C]{\small\color{muted}\thepage}
\renewcommand{\headrulewidth}{0.3pt}

\newcommand{\code}[1]{\texttt{#1}}
\newcommand{\model}[1]{\textbf{#1}}
\newcommand{\good}[1]{\textcolor{green}{\textbf{#1}}}
\newcommand{\warn}[1]{\textcolor{orange}{\textbf{#1}}}
\newcommand{\reportbox}[2]{%
  \begin{center}
  \fcolorbox{#1!45}{#1!6}{%
    \begin{minipage}{0.94\linewidth}
      \small #2
    \end{minipage}}
  \end{center}}

\title{\vspace{-1.0cm}
  \color{navy}\bfseries Cross-model transfer of evolved Higgs FeatureGraphs\\[4pt]
  \large CatBoost and TabM, five matched estimator seeds\\
  \normalsize Higgs-small binary classification}
\author{GigaEvo research report}
\date{23 July 2026}

\begin{document}
\maketitle
\vspace{-0.7cm}

\begin{abstract}
\noindent
We evolved one FeatureGraph against CatBoost and one against a paper-tuned
Higgs TabM-mini$^\dagger$, froze the CV-selected winners, and evaluated both
graphs with both estimators on the same untouched test split.  The resulting
two-by-two transfer matrix, plus raw-feature controls, contains 30 fresh model
fits: five matched seeds for each of two evaluators and three feature inputs.
Every engineered cell beats its row's raw control on every seed.  More
surprisingly, the TabM-evolved graph is best for both estimators.  It raises
CatBoost accuracy from $72.728\%\pm0.083\%$ to
$74.010\%\pm0.071\%$ and TabM accuracy from
$73.770\%\pm0.113\%$ to $74.871\%\pm0.161\%$.
The two searches emit no identically named feature, yet code inspection and
training-row Spearman correlations recover four shared mechanisms.  CatBoost
builds broad event summaries, mass ratios, angular distances, and a b-tag
aggregate.  TabM builds a denser and more internally consistent geometric
representation: scaled angular distances, momentum closure, longitudinal
centrality, and lepton--jet mass proxies.  That geometry is not specific to a
neural consumer; it transfers especially well to trees.
\end{abstract}

\reportbox{blue}{
\textbf{\color{navy}Bottom line.}
The TabM search discovered the stronger \emph{feature program}, not merely a
TabM-specific adapter.  Its graph improves CatBoost test accuracy by
\good{+1.282 percentage points}, versus +0.797 for CatBoost's own graph, and
even lets CatBoost with evolved features exceed raw TabM by 0.241 points.
The tuned TabM plus its own graph remains the overall winner at 74.871\%.
This is one fixed split and one evolutionary run per estimator, so the
transfer pattern is a strong hypothesis for replication rather than an
estimate of dataset-level uncertainty.}

\begin{figure}[H]
  \centering
  \includegraphics[width=\linewidth]{figures/performance_matrix.pdf}
  \caption{\textbf{Held-out performance.} Rows are freshly fitted evaluators;
  columns are frozen feature programs.  Values are mean $\pm$ sample SD over
  seeds 0--4.  Orange outlines mark row winners.}
  \label{fig:performance}
\end{figure}

\section{Question and protocol}

\subsection{What is transferred}

A FeatureGraph is executable feature code, not a fitted prediction model.
The two JSON graphs were selected by maximum valid three-fold CV accuracy
during separate 100-mutation Memory V2 campaigns.  Selection never read test
labels.  We then froze the exact JSON bytes and evaluated:

\[
\{\text{CatBoost evaluator},\ \text{TabM evaluator}\}
\ \times\
\{\text{raw},\ \text{CatBoost graph},\ \text{TabM graph}\}.
\]

Each cell rebuilds graph preprocessing and the downstream estimator.  For
CatBoost and TabM, validation selects the training length, the search model is
discarded, and a fresh estimator is trained on train plus validation for that
length before predicting test.  Only the estimator seed changes across the
five repetitions; the graph, train/validation/test split, and test rows stay
fixed.  Consequently, the reported SD measures training randomness, not
sampling uncertainty of Higgs events.

Higgs-small contains 62,751 training, 15,688 validation, and 19,610 test
events, with 28 numerical inputs.  The first 21 columns describe lepton,
missing-energy, and four-jet kinematics (including four discrete b-tag
discriminants); the final seven are derived invariant-mass variables.

\reportbox{orange}{
\textbf{Model-recipe boundary.}
This matrix deliberately uses the same paper-tuned Higgs
TabM-mini$^\dagger$ recipe used during evolution: $k=32$, three 816-wide
blocks, dropout 0.4325, 7-bin 20-dimensional PLE, AdamW learning rate
$9.498\times10^{-4}$, batch 512, and BF16.  It is a labeled reproduction
setting.  Future ordinary baselines use the repository's single
dataset-independent example recipe---full TabM, $k=32$, $2\times512$,
dropout 0.1, 48-bin 16-dimensional PLE, AdamW 0.002/0.0003---without
per-dataset architecture selection.}

\subsection{Evolutionary selection}

\begin{table}[H]
\centering
\small
\begin{tabular}{lrrrrr}
\toprule
Search evaluator & Raw CV accuracy & Winner CV accuracy & Gain (pp) &
Iteration & Outputs \\
\midrule
\input{results/evolution_rows.tex}
\end{tabular}
\caption{CV selection before test access.  Dispersion is the three-fold CV
SD stored with each program.}
\label{tab:evolution}
\end{table}

Both searches improve their native CV evaluator by roughly 0.9 points.
CatBoost's winner has ten nodes, ten outputs, and maximum graph depth two.
TabM's winner has only five independent nodes but emits fifteen outputs at
depth one.  All original features remain available in both graphs.

\section{Cross-transfer results}

\begin{table}[H]
\centering
\small
\begin{tabular}{llr@{\hspace{1.2em}}rr@{\hspace{1.2em}}r}
\toprule
Evaluator & Graph & Accuracy (\%) & Gain (pp) & ROC AUC (\%) & Gain (pp) \\
\midrule
\input{results/test_rows.tex}
\end{tabular}
\caption{Five-seed test means $\pm$ sample SD.  Gain is the mean matched-seed
difference from the raw control in the same evaluator row.}
\label{tab:test}
\end{table}

\begin{figure}[H]
  \centering
  \includegraphics[width=0.88\linewidth]{figures/paired_gain_matrix.pdf}
  \caption{\textbf{Paired uplift over raw inputs.}  Positive values occur on
  every individual seed, not only after averaging.}
  \label{fig:uplift}
\end{figure}

Three comparisons matter:

\begin{itemize}
  \item \model{TabM's graph wins the CatBoost row.}  Its +1.282-point
  accuracy gain exceeds CatBoost's native +0.797 by
  $0.485\pm0.099$ points across matched seeds.  The analogous ROC-AUC
  advantage is $0.495\pm0.118$ points.
  \item \model{TabM also strongly prefers its own graph.}  Its native gain is
  +1.101 accuracy points, while the CatBoost graph transfers positively but
  yields only +0.411.  The native-over-foreign difference is
  $0.690\pm0.205$ points.
  \item \model{Features narrow but do not erase the model gap.}  Raw TabM is
  1.041 accuracy points ahead of raw CatBoost.  Under the common TabM graph
  the gap is 0.861 points.  CatBoost plus the TabM graph does, however,
  exceed raw TabM by 0.241 points.
\end{itemize}

These are paired descriptive contrasts: seed 0 is subtracted from seed 0,
and so on.  They are more precise than comparing two independent SDs because
the held-out rows are shared, but five estimator seeds do not quantify
variation from rerunning evolution or resampling the dataset.

\section{What the two searches mined}

\begin{figure}[H]
  \centering
  \includegraphics[width=\linewidth]{figures/feature_anatomy.pdf}
  \caption{\textbf{Feature semantics and graph structure.} Each output is
  assigned to exactly one semantic family by manual code inspection.  Counts
  describe emitted columns, not learned importance.}
  \label{fig:anatomy}
\end{figure}

\subsection{CatBoost: broad, one-concept-per-node coverage}

CatBoost emits ten scalar features from ten nodes:

\begin{itemize}
  \item three angular separations:
  $\Delta R(\ell,j_1)$, $\Delta R(\ell,j_2)$, and $\Delta R(j_1,j_2)$;
  \item a lepton--MET transverse-mass proxy and a lepton--MET transverse
  vector-sum magnitude;
  \item an event-scale sum, leading-jet fraction, and MET-to-scale ratio;
  \item a shifted $m_{bb}/m_{jj}$ ratio and the sum of four b-tag
  discriminants.
\end{itemize}

This graph touches 19 raw columns and uniquely uses both the supplied
high-level mass features and flavor information.  Two ratio nodes depend on
the event-scale node, creating depth two.  Its coverage is deliberately broad:
angular geometry, energy allocation, mass structure, and flavor each receive
at least one explicit coordinate.

\subsection{TabM: compact nodes, dense event geometry}

TabM emits fifteen features from five multi-output nodes:

\begin{itemize}
  \item $\Delta R(\ell,j_1)$, $\Delta R(j_1,j_2)$, a MET/visible-$p_T$
  ratio, and lepton--MET transverse mass;
  \item visible transverse-vector magnitude, residual momentum imbalance,
  and visible--MET angular alignment;
  \item six Zeppenfeld-like coordinates: normalized longitudinal position and
  a bounded centrality score for the lepton, jet 3, and jet 4 relative to the
  leading-jet pair;
  \item two lepton--jet mass proxies.
\end{itemize}

It touches 17 raw columns, almost entirely low-level kinematics.  It does not
engineer b-tags or reuse the seven supplied mass columns.  Instead, six of
its fifteen outputs describe longitudinal centrality and three describe
transverse momentum closure.  This is a coherent geometric basis rather than
a broad list of independent event summaries.

Both winners are entirely rowwise and target-free: there are no learned
target encodings, aggregate nodes, dropped raw columns, or target
transformations.  Their gains cannot be attributed to supervised
feature-state leakage.

\section{How similar are the feature sets?}

Exact output-name Jaccard similarity is zero: the ten CatBoost names and
fifteen TabM names have no literal intersection.  That lexical view is
misleading.  Formula inspection identifies four shared mechanisms:
lepton--jet-1 distance, jet-1--jet-2 distance, lepton--MET transverse mass,
and missing energy normalized by an event scale.

\begin{figure}[H]
  \centering
  \includegraphics[width=\linewidth]{figures/cross_feature_correlation.pdf}
  \caption{\textbf{Cross-graph Spearman correlation on all 62,751 training
  rows.} Only cells with $|\rho|\geq0.60$ are annotated.  Correlation detects
  functional similarity but is not a feature-importance measure.}
  \label{fig:correlation}
\end{figure}

\begin{table}[H]
\centering
\small
\begin{tabular}{llr}
\toprule
CatBoost output & TabM output & Spearman $\rho$ \\
\midrule
\input{results/top_correlations_rows.tex}
\end{tabular}
\caption{Seven largest absolute cross-graph correlations.}
\label{tab:correlations}
\end{table}

Two shared distances are exactly rank-identical ($\rho=1$), despite different
names.  The two event-scale ratios are nearly equivalent
($\rho=0.993$).  The transverse-mass features are related but not identical
($\rho=0.797$): TabM consistently rescales normalized azimuthal coordinates
by 1.8138 before trigonometry, whereas CatBoost's transverse-mass node does
not.  Lepton--jet distance also correlates strongly with the corresponding
TabM mass proxy because both contain the same angular separation.

The low-correlation regions are at least as informative.  CatBoost's b-tag
sum and high-level mass ratio have no close TabM counterpart.  Conversely,
TabM's signed longitudinal positions $z_\ell,z_{j_3},z_{j_4}$ are nearly
orthogonal to every CatBoost output; the associated bounded centralities and
full momentum-closure features form a genuinely new block of representation.

\section{Interpretation, limitations, and reproducibility}

\subsection{Why the TabM graph plausibly transfers to trees}

The TabM graph presents nonlinear interactions as explicit scalar
coordinates.  A depth-six tree no longer has to approximate circular
differences, vector addition, ratios, and centrality from many axis-aligned
splits.  The same representation also helps a shallow neural ensemble.
Consistent angular rescaling may explain part of the advantage over the
CatBoost graph, while the six centrality outputs and three momentum-closure
outputs supply information absent from it.  This is a mechanism-level
interpretation, not a causal ablation.

The names ``transverse mass'' and ``invariant mass'' require care.  The
Higgs-small inputs are normalized; azimuthal ranges show that multiplying by
1.8138 approximately restores radians, but momentum and mass columns are not
in GeV.  The engineered quantities are therefore physics-inspired nonlinear
proxies, not calibrated physical observables.  CatBoost's $+3$ shifts in its
event and mass ratios are likewise numerical heuristics rather than unit
restoration.

\subsection{Limits}

\begin{itemize}
  \item One evolution run per estimator cannot separate stable search behavior
  from winner-selection luck.
  \item All 30 fits reuse one official test split.  Seed SD excludes split,
  finite-sample, and simulation-distribution uncertainty.
  \item The matrix is a post-hoc audit.  Choosing the TabM graph because it
  wins this test matrix would spend the test set; future datasets should
  pre-register the donor or select it using validation only.
  \item TabM emits five more columns than CatBoost.  No size-matched or
  leave-one-family-out ablation was run, so geometry and feature count are
  partially confounded.
  \item The campaigns used Memory V2.  This report studies the final feature
  programs; it does not estimate a causal effect of memory.
\end{itemize}

\subsection{Artifact trail}

The report directory contains the exact graph JSON files, their SHA-256
digests, CV-selection provenance, all 30 seed-level test records, derived CSV
tables, and scripts for freezing, evaluation, plotting, and compilation.
\code{run\_cross\_eval.sh} pins the historical Higgs TabM recipe explicitly;
normal TabM runs do not inherit it.  Reproduction commands and file roles are
listed in \code{README.md} and \code{MANIFEST.md}.

\end{document}
